\begin{question}[index=17,points=10]
设 $\Gamma$ 为 $x^2+y^2=2x$（$y\ge0$）上从 $O(0,0)$ 到 $A(2,0)$ 的一段弧，连续函数 $f(x)$ 满足
\[f(x)=x^2+\int_\Gamma y[f(x)+\eu^x]\dd x+(\eu^x-xy^2)\dd y,\]
求 $f(x)$.
\end{question}
\QuestionBreak\WorkSpace

\begin{question}[index=18,points=12]
求 $f(x,y)=\eu^{x+y}\left(y+\dfrac{x^3}3\right)$ 的极值.
\end{question}
\QuestionBreak\WorkSpace\WorkPageBreak

\begin{question}[index=19,points=12]
已知函数 $f(x,y)$ 二阶偏导数连续，任给 $x,y$，均有 $f(x+1,y)=f(x,y)$，$f(x,y+1)=f(x,y)$，且
\[\iint_D\{[f'_x(x,y)]^2+[f'_y(x,y)]^2\}\dd\sigma=a,\quad D=\{(x,y)\mid-1\le x\le1,-1\le y\le1\}.\]
\par（1）计算 $\displaystyle I=\iint_Df(x,y)[f''_{xx}(x,y)+f''_{yy}(x,y)]\dd\sigma$；
\par（2）若 $I\ge0$，证明 $\mathrm d[f(x,y)]=0$.
\end{question}
\QuestionBreak\WorkSpace

\begin{question}[index=20,points=12]
设 $a,b$ 为实数，函数 $z=1+ax^2+by^2$ 在点 $(1,1)$ 处的方向导数中，沿方向 $\vect l=2\vect i+4\vect j$ 的方向导数最大，且最大值为 $2\sqrt5$.
\par（1）求 $a,b$；
\par（2）求曲面 $z=1+ax^2+by^2$ 被曲面 $z=2(x^2+3y^2)$ 所截部分的面积.
\end{question}
\QuestionBreak\WorkSpace\WorkPageBreak

\begin{question}[index=21,points=12]
设矩阵 $\vect A,\vect B$ 满足：对任意 $x_1,x_2,x_3$，均有
\[\vect A\mat{x_1\\x_2}=\mat{x_1+2x_2\\x_1-x_2\\2x_1+ax_2},\qquad\vect B\mat{x_1\\x_2\\x_3}=\mat{x_1+2x_2+3x_3\\bx_1-2x_2-3x_3}.\]
若 $\vect A^{\mathrm T}\vect x=\vect0$ 的解均为 $\vect B\vect x=\vect0$ 的解且它们不同解.
\par（1）求 $a,b$ 的值；
\par（2）用正交矩阵 $\vect Q$ 将 $f(x_1,x_2)=\vect x^{\mathrm T}\vect B\vect A\vect x$ 化为标准形.
\end{question}
\QuestionBreak\WorkSpace

\begin{question}[index=22,points=12]
某通信网络中，数据传输过程需经过 $3$ 个独立的中间结点，每个结点的处理延迟时间（单位：秒）服从参数为 $\lambda$ 的指数分布，平均延迟为 $0.2$ 秒，总延迟时间 $X$ 为 $3$ 个结点延迟时间之和.
\par（1）求 $X$ 的概率密度；
\par（2）计算 $P\{X<0.5\}$.
\end{question}
\WorkSpace
